\section{Methods}
\label{sec:methods}

\subsection{Dataset and preprocessing}

We used the CPTAC-LUAD cohort described by Gillette et al.\ \cite{gillette2020proteogenomic},
accessed via the \texttt{cptac} Python package.
Of the 111 patients in the dataset, nine were excluded because they lacked a complete
matched tumour-normal pair in at least one of the three omics modalities, leaving 102 patients
(204 samples) for analysis.
The three modalities were RNA-seq (BCM pipeline, reported as TPM and converted to
$\log_2(\text{TPM}+1)$), proteomics (Michigan CPTAC pipeline, log-ratio), and
phosphoproteomics (Michigan, log-ratio).
Missing values were handled as follows: features absent in more than 20\% of samples were
removed; the remainder were median-imputed column-wise.

Feature selection was applied independently per modality using the median absolute
deviation (MAD) of each feature across all 204 samples.
The RNA-seq data contained 59,286 features, reflecting the comprehensive GENCODE annotation
used by the BCM pipeline and including protein-coding genes, long non-coding RNAs,
and pseudogenes.
We retained the top 5,000 transcripts, 5,000 proteins, and 6,000 phosphosites by MAD.
The retained features were then z-scored per feature using a single scaler fit jointly on
all 204 samples; tumour and normal samples were not z-scored separately.
This joint scaling is required for the $\Delta Z$ framework: both tissue states must share
a common reference frame for their difference to be interpretable as a biological
transformation.

\subsection{Latent factor modelling with MOFA+}

We applied Multi-Omics Factor Analysis (MOFA+, \cite{argelaguet2020mofa2}) to all 204
samples treated as a single group with three views.
MOFA+ decomposes the observed data matrix as a product of per-sample latent factor scores
and per-feature loading weights, assuming that each observed value arises from a
Gaussian distribution centred on the linear combination of factor contributions.

We initialised 20 factors and fit five models from different random seeds
(convergence mode: fast, maximum 1,000 iterations); the model with the highest
final evidence lower bound (ELBO) was selected for downstream analysis.

Inactive factors were removed post-hoc: a factor was retained only if the fraction of
variance it explained in at least one view exceeded a threshold of $R^2_{\max} \geq 0.01$.
This retained 13 of the initial 20 factors.

\subsection{Transformation vectors and clustering}

For each of the 102 patients we computed a latent transformation vector as defined
in equation~(\ref{eq:deltaz}), using the 13-dimensional factor score vector of each sample.
This produces a $102 \times 13$ matrix whose rows encode the direction and magnitude
of each patient's tumourigenesis in the shared MOFA latent space.

In matched tumour-normal data, the first MOFA factor typically captures the dominant
axis of variation, which is the tissue-type contrast itself.
Variation along this axis reflects the overall extent of transcriptional reprogramming
and is shared across patients; it therefore carries less information about
patient-specific transformation patterns than the remaining factors.
We constructed four representations of the $\Delta Z$ matrix to assess sensitivity
to this: raw $\Delta Z$ using all 13 factors; raw $\Delta Z$ with Factor~1 excluded;
L2-normalised $\Delta Z$ (all factors); and L2-normalised $\Delta Z$ with Factor~1
excluded.
L2 normalisation places all patients on a unit hypersphere, making clustering
respond to transformation direction rather than magnitude.
The primary representation for clustering was L2-normalised $\Delta Z$ on all 13 factors.

We clustered the 102 patients using $k$-means (\texttt{sklearn.cluster.KMeans},
\texttt{n\_init=100}), sweeping $k \in \{2, 3, 4, 5, 6\}$ and recording the mean
silhouette score at each $k$.
Gaussian Mixture Model (GMM) AIC and BIC were computed at each $k$ as additional
descriptive metrics.
We selected $k = 3$ based on the silhouette profile, cluster balance (37, 25, and
40 patients), and partition robustness across representations.
Robustness was quantified using the adjusted Rand index (ARI), which measures
agreement between two independent cluster assignments on a scale from 0 (chance)
to 1 (identical).
Re-fitting the same algorithm on the three alternative representations and comparing
with the primary partition yielded ARI $= 0.82$ against L2-normalised $\Delta Z$
without Factor~1, and ARI $= 0.80$ against raw $\Delta Z$, indicating that the
three-subtype structure is stable across representational choices.
Silhouette scores in the range $0.18$--$0.19$ are lower than those typical of
well-separated synthetic clusters; however, values in this range are biologically
plausible for continuous molecular phenotypes and the decision to proceed with
$k = 3$ was supported by the consistency of the partition across representations.

\subsection{Biological and clinical interpretation}

\paragraph{Factor interpretation.}
MOFA factor loadings measure the contribution of each gene to a given factor: genes
with large positive loadings are the primary drivers of high factor scores, and genes
with large negative loadings drive low scores.
To characterise the biological programmes encoded in each factor, we ranked all 5,000
transcripts by their RNA loading weight and applied preranked gene set enrichment
analysis (GSEA).
Preranked GSEA walks down a sorted gene list and tests whether the members of each
predefined gene set are concentrated near the top or bottom relative to a permutation
null; the normalised enrichment score (NES) indicates the direction and strength of
enrichment.
A positive NES indicates enrichment among genes with large positive loadings; a
negative NES indicates enrichment at the negative-loading pole.
We used the MSigDB Hallmark 2020 and KEGG 2021 Human libraries (1,000 permutations,
gene set size 15--500) and filtered results at FDR $< 0.10$ (Benjamini-Hochberg).

\paragraph{Cluster interpretation.}
To test whether the transformation subtypes differ in their transcriptional
reprogramming, we computed per-patient differential expression as
$\Delta \text{Expr}_i = \text{RNA}_{\text{tumour},i} - \text{RNA}_{\text{normal},i}$
across the full 59,286-gene matrix from ingestion rather than the MAD-filtered
MOFA input.
This avoids biasing the test towards features that were pre-selected for latent
factor construction; the cluster interpretation is thereby an independent assessment.
For each subtype we performed one-vs-rest Mann-Whitney $U$ tests on
$\Delta \text{Expr}$ across all genes and applied Benjamini-Hochberg correction.
The per-gene median difference between the focal cluster and all remaining patients
was used as the ranking statistic for a second round of preranked GSEA with the
same libraries and FDR threshold.

\paragraph{Clinical associations.}
Overall survival was estimated using the Kaplan-Meier method for each transformation
subtype.
An omnibus log-rank test was applied to assess whether survival distributions differed
across the three subtypes simultaneously.
Pairwise log-rank comparisons were conducted only when the omnibus test yielded
$p < 0.05$, with Benjamini-Hochberg correction applied.
Enrichment of driver mutations (KRAS, EGFR, TP53, ALK) within each subtype was
evaluated by two-sided Fisher's exact test, corrected across all mutation-subtype
pairs by Benjamini-Hochberg.
Given $n = 102$ patients, all survival analyses should be interpreted as
hypothesis-generating; the small cohort size makes it this analysis severely underpowered.
Full configuration parameters are available in the project repository at
\texttt{config/config.yml}.
